\chapter{Ball Volumes and the Gamma Connection}\label{ch:n-ball-gamma}

\section{An exact represented formula}
Write
\[
 M_n(p,R)=c_np^{\lfloor n/2\rfloor}R^n,\qquad
 c_0=1,\quad c_1=2,\quad c_{n+2}=\frac{2c_n}{n+2}.
\]
The existing finite volume model proves its recurrence and enclosure bounds.
The new endpoint evaluator extends this computation to supplied valid
representations of arbitrary nonnegative inputs, including irrational radii.
Its hypotheses require nonnegative lower endpoints at every stage.

\begin{theorem}[Represented ball formula]
\label{thm:n-ball-represented}
The endpoint evaluation of $M_n$ is valid for every natural dimension and
every pair of valid nonnegative interval presentations. Replacing either
presentation by an equivalent valid nonnegative presentation preserves its
value under \texttt{RealRaw.Equiv}.
\lean{ComputableAnalysis.NBallRaw.volume_valid, ComputableAnalysis.NBallRaw.volume_equiv}
\leanok
\end{theorem}

This is a theorem about the formula. Identifying that formula with geometric
ball volume remains a separate task. The validity proof inherits a native
computation of the fixed rational fact $0<2$ from the existing multiplication
foundation; the publication records that dependency explicitly.

\begin{theorem}[Half-step coefficient cancellation]
\label{thm:n-ball-gamma-coefficients}
Define $g_0=1$, $g_1=1/2$, and $g_{n+2}=(n+2)g_n/2$. Then
$c_ng_n=1$ and $M_n(p,R)g_n=p^{\lfloor n/2\rfloor}R^n$.
\lean{ComputableAnalysis.nBallCoeff_mul_gammaHalfCoeff, ComputableAnalysis.nBallVolumeModel_gamma}
\leanok
\end{theorem}

\section{Finite progress toward Gaussian normalization}
\begin{theorem}[Square, triangles, and diagonal]
\label{thm:gaussian-square-diagonal}
For any finite list of rational cell masses,
\[
 (\sum_i a_i)^2=2\sum_{i<j}a_ia_j+\sum_i a_i^2.
\]
If $0\le a_i\le\delta$, then $\sum_i a_i^2\le\delta\sum_i a_i$.
\lean{ComputableAnalysis.gaussian_square_split, ComputableAnalysis.gaussian_diagonal_bound}
\leanok
\end{theorem}
The diagonal estimate supports the bounded slope substitution described in
Chapter~\ref{ch:gaussian-convolution}; it does not prove that substitution.

\section{Geometric and analytic targets}
Construct the Gamma integral and its recurrence with explicit endpoint and
tail estimates. Normalize the Gaussian independently against geometric $\pi$.
Then compare Cartesian products with radial shells to justify
\[
 \pi^{n/2}\simeq nV_n(1)\int_0^\infty e^{-r^2}r^{n-1}\,dr
 \simeq V_n(1)\Gamma(n/2+1),\qquad n\ge1.
\]
This should establish $V_n(R)\simeq\pi^{n/2}R^n/\Gamma(n/2+1)$.
The independent coefficient cancellation above is not yet this integral or
geometric theorem. The separate reader page supplies a rational interval
calculator, tested against Lean evaluations, and explains these remaining
bridges. No Mathlib real numbers enter the construction.
